% IEEE T-IV Journal paper — IEEEtran two-column format
\documentclass[journal]{IEEEtran}

% ——— packages ———
\usepackage[T1]{fontenc}
\usepackage[utf8]{inputenc}
\usepackage{amsmath,amssymb,amsfonts}
\usepackage{graphicx}
\usepackage{booktabs}
\usepackage{multirow}
\usepackage{cite}
\usepackage{url}
\usepackage{hyperref}
\usepackage{xcolor}
\usepackage{textcomp}
\usepackage{balance}
\usepackage{adjustbox}

% figure path
\graphicspath{{./}}

% tilde in text
\newcommand{\tsim}{{\raise.17ex\hbox{$\scriptstyle\sim$}}}

% ——— document ———
\begin{document}

\title{Quantifying the Relative Impact of Class Imbalance Handling and Domain Configuration on Vehicle-Dynamics-Based Drowsy Driving Detection}

\author{%
  \IEEEauthorblockN{First~Author\IEEEauthorrefmark{1},
                     Second~Author\IEEEauthorrefmark{2}}\\
  \IEEEauthorblockA{\IEEEauthorrefmark{1}Department of XXXX, University of XXXX, City, Country\\
  \IEEEauthorrefmark{2}Department of XXXX, University of XXXX, City, Country\\
  Email: \{first.author, second.author\}@example.ac.jp}
}

\maketitle

% ---------- ABSTRACT ----------
\begin{abstract}
Drowsy driving detection (DDD) using vehicle dynamics faces two intertwined challenges: class imbalance and inter-driver domain shift.
Existing studies address these in isolation, leaving practitioners without guidance on which design choice matters most.
We introduce a factorial framework quantifying the relative importance of four factors---rebalancing strategy ($R$), training mode ($M$), distance metric ($D$), and domain membership ($G$)---via Sobol--Hoeffding variance decomposition over 1,512 random forest evaluations on an 87-driver public simulator dataset.
Training mode and rebalancing jointly account for ${>}95$\% of systematic variance, with a substantial $R \times M$ interaction causing strategy ranking reversal across modes ($\rho = -0.74$ to $-0.89$).
Distance metric and domain membership are negligible.
SW-SMOTE at ratio 0.1 achieves the best within-domain F2 (0.558) and AUROC (0.903).
A full replication with XGBoost (1,512 additional evaluations) confirms that the factor hierarchy $S_{TM} > S_{TR} \gg S_{TD} \approx S_{TG}$ is classifier-independent.
Vehicle dynamics coupling provides a physics-grounded explanation for domain-configuration irrelevance.
Practitioners should prioritize training mode and class rebalancing over domain grouping.
\end{abstract}

\begin{IEEEkeywords}
Drowsy driving detection, class imbalance, domain shift, vehicle dynamics, Sobol sensitivity analysis, factorial experiment, SMOTE.
\end{IEEEkeywords}

% =====================================================================
\section{Introduction}
\label{sec:introduction}

Drowsy driving is a leading cause of traffic accidents worldwide, accounting for an estimated 15--20\% of road fatalities~\cite{who2018}.
Vehicle-dynamics-based drowsy driving detection (DDD) is an attractive non-intrusive approach that monitors signals such as steering angle, lateral acceleration, and lane offset without requiring sensors attached to the driver.
However, deploying such systems in practice requires addressing two well-documented challenges that degrade classifier performance.

First, \textbf{class imbalance}: in naturalistic driving data, alert states vastly outnumber drowsy states (typically 9:1 or greater), biasing classifiers toward the majority class and degrading detection of safety-critical drowsy episodes.
Resampling techniques such as SMOTE~\cite{chawla2002} and random undersampling (RUS)~\cite{he2009} are commonly applied, but their relative effectiveness for DDD has not been systematically evaluated across operational conditions.

Second, \textbf{domain shift}: inter-subject variability in driving behavior causes performance degradation when a model trained on one group of drivers is applied to another~\cite{pan2010}.
Practitioners must choose a distance metric to quantify subject dissimilarity, determine domain membership, and select a training mode (cross-domain, within-domain, or mixed).
Each of these decisions constitutes a design factor, yet their relative contribution to downstream classification performance remains unknown.

Crucially, these two challenges are not independent.
A rebalancing strategy that improves performance under one training mode may degrade it under another.
Yet prior work has studied class imbalance handling~\cite{chawla2002,he2009} and domain adaptation~\cite{pan2010,weiss2016} in isolation, treating them as separate optimization targets rather than as interacting components of a unified experimental design.
No study has quantified their relative importance or characterized their interactions.

This paper addresses this gap by framing DDD system design as a \textbf{four-factor factorial experiment} and applying variance-based sensitivity analysis to decompose performance variance into attributable sources.
Specifically, we:
\begin{enumerate}
  \item Conduct an intuitive one-factor-at-a-time (OFAT) analysis to visualize each factor's marginal effect across all fixed conditions, identifying which factors produce consistent performance shifts and which are absorbed by other design choices.
  \item Apply Sobol--Hoeffding variance decomposition to quantify each factor's first-order and total-order contributions, revealing that training mode ($M$) and rebalancing strategy ($R$) jointly account for ${>}\,95$\% of systematic variance, while distance metric ($D$) and domain membership ($G$) are negligible.
  \item Characterize the $R \times M$ interaction---the third-largest effect---which produces a complete strategy ranking reversal between cross-domain and within-domain settings, demonstrating that these two factors cannot be optimized independently.
  \item Provide a physics-grounded explanation rooted in the bicycle model for why distance metric choice is irrelevant despite generating genuinely different subject rankings.
\end{enumerate}

To the best of our knowledge, this is the first study to apply a factorial experimental design with Sobol sensitivity analysis to DDD, systematically decomposing the relative contribution of class imbalance handling and domain configuration decisions.

% =====================================================================
\section{Related Work}
\label{sec:related}

\subsection{Vehicle-Dynamics-Based Drowsy Driving Detection}

Vehicle dynamics signals---steering wheel angle, lateral acceleration, and lane position---are widely used for DDD owing to their non-intrusive nature.
Arefnezhad et~al.~\cite{arefnezhad2019} used steering wheel angle and velocity with statistical and spectral features for drowsiness classification.
Wang et~al.~\cite{wang2022} extended this approach using LSTM networks with five vehicle signals, applying Gaussian smoothing with a 1-second window.
Zhao et~al.~\cite{zhao2012} applied continuous wavelet transform to capture transient changes in driving behavior.
Recent surveys~\cite{fonseca2025,alquraishi2024} confirm that vehicle-dynamics-based methods remain a practical backbone, but note that most studies evaluate a single architecture on a fixed data split without systematically examining the effect of preprocessing decisions.

\subsection{Class Imbalance in Safety-Critical Detection}

Class imbalance is inherent in drowsy driving data: drowsy episodes are rare but safety-critical.
He and Garcia~\cite{he2009} provide a comprehensive taxonomy of methods, including random undersampling (RUS), SMOTE~\cite{chawla2002}, and cost-sensitive learning.
Beyond standard SMOTE, variants such as Borderline-SMOTE~\cite{han2005} and ADASYN~\cite{he2008adasyn} adapt synthesis to decision-boundary regions; however, for multi-subject pooled data, subject-wise SMOTE (SW-SMOTE)---which generates synthetic samples within each subject's feature subspace before pooling---better preserves individual driving patterns and avoids cross-subject artifacts.
For DDD, He et~al.~\cite{he2024} proposed a generative adaptive CNN to address class-imbalanced EEG data.
However, most DDD studies either ignore imbalance entirely or apply a single method without comparing alternatives.
No study has evaluated the interaction between rebalancing strategy and training mode in a factorial framework.

\subsection{Domain Shift and Cross-Subject Generalization}

Inter-subject variability creates domain shift in DDD~\cite{pan2010,aravinth2025}.
Kim et~al.~\cite{kim2025} proposed calibration-free drowsiness classification with prototype-based domain mixup for EEG signals.
Luo et~al.~\cite{luo2024} developed cross-scenario and cross-subject domain adaptation methods.
Li et~al.~\cite{li2026} introduced MADNet for EEG-based cross-subject drowsiness diagnosis using domain generalization.
These studies focus on architectural innovations to mitigate domain shift but do not quantify the relative importance of domain-related design choices (distance metric, membership, training mode) compared to preprocessing choices (rebalancing).

\subsection{Sensitivity Analysis in Machine Learning}

Variance-based sensitivity analysis using Sobol indices~\cite{sobol1993} has been applied in engineering and biological modeling~\cite{tosin2020} but remains underutilized in machine learning.
Taylor et~al.~\cite{taylor2021} pioneered the application of Sobol indices to rank deep learning hyperparameter influence, demonstrating that learning rate and batch size dominate while architectural choices have lower impact.
Kuhnt and Kalka~\cite{kuhnt2022} used Sobol indices with FANOVA graphs for ML model interpretation.
These works demonstrate the power of Sobol decomposition for identifying dominant factors in complex systems, but no study has applied this methodology to DDD or to the interplay between data preprocessing and domain configuration.

% =====================================================================
\section{Methodology}
\label{sec:methodology}

The proposed methodology proceeds in five stages: (1)~data acquisition and feature engineering from vehicle dynamics; (2)~domain grouping of subjects via pairwise distance metrics; (3)~factorial enumeration of all design factor combinations; (4)~independent RF training with Optuna hyperparameter optimization for each cell; and (5)~Sobol--Hoeffding variance decomposition to rank factor importance.
Stages~1--4 produce a balanced observation matrix of 1,512 classification scores; stage~5 decomposes the variance of this matrix into attributable sources.
Fig.~\ref{fig:pipeline} provides a system overview.

\begin{figure}[!t]
  \centering
  \includegraphics[width=\columnwidth]{fig1_pipeline_overview.pdf}
  \caption{System overview of the proposed factorial sensitivity analysis pipeline for vehicle-dynamics-based DDD. Raw simulator signals are featurized, subjects are grouped by pairwise distance, and all 126 factor combinations are evaluated over 12 seeds. The resulting 1,512 observations are decomposed via Sobol--Hoeffding analysis.}
  \label{fig:pipeline}
\end{figure}

\subsection{Data Acquisition and Feature Extraction}
\label{sec:data}

The publicly available multi-modal driving dataset of Aygun et~al.~\cite{aygun2024} (Harvard Dataverse, DOI: 10.7910/DVN/HMZ5RG) is used, comprising 87 subjects with two sessions each recorded on the SIMlsl driving simulator at $f_s = 60$\,Hz.
Five raw signals are extracted (Table~\ref{tab:signals}).

\begin{table}[!t]
  \caption{Vehicle Dynamics Signals}
  \label{tab:signals}
  \centering
  \begin{tabular}{c l c}
    \toprule
    Symbol & Signal & Unit \\
    \midrule
    $\delta(t)$              & Steering angle           & rad \\
    $\dot{\delta}(t)$        & Steering speed           & rad/s \\
    $a_y(t)$                 & Lateral acceleration     & m/s$^2$ \\
    $a_x(t)$                 & Longitudinal acceleration & m/s$^2$ \\
    $e_{\mathrm{lane}}(t)$   & Lane offset              & m \\
    \bottomrule
  \end{tabular}
\end{table}

Under the linear bicycle model~\cite{rajamani2012}, lateral dynamics follow:
\begin{equation}
  a_y = v \cdot \frac{\dot{\delta}}{L}
  \label{eq:bicycle}
\end{equation}
where $v$ is vehicle speed and $L$ is the wheelbase.
Lane offset evolves as:
\begin{equation}
  \ddot{e}_{\mathrm{lane}}(t) = a_y(t) - a_{y,\mathrm{road}}(t)
  \label{eq:lane}
\end{equation}
creating a four-signal causal chain $\delta \to \dot{\delta} \to a_y \to e_{\mathrm{lane}}$, with only $a_x$ dynamically independent.

Raw signals are segmented into 3-second sliding windows with 50\% overlap (1.5-second step).
Each window is labeled using the Karolinska Sleepiness Scale (KSS)~\cite{akerstedt1990}: KSS\,$\in\{1,\ldots,5\}$ maps to Alert (class~0) and KSS\,$\in\{8,9\}$ to Drowsy (class~1); intermediate scores (6,\,7) are excluded to maximize label reliability.

A total of 135 features are extracted per window using four methods:
(i)~statistical and spectral features ($22\times 2 = 44$ dimensions) from FFT power spectra in the $[0.5,\,30]$\,Hz band;
(ii)~smooth/std/prediction-error features ($3\times 5 = 15$ dimensions) using 2nd-order Taylor approximation~\cite{atiquzzaman2018};
(iii)~permutation entropy features~\cite{bandt2002} ($8\times 5 = 40$ dimensions) capturing signal complexity; and
(iv)~time-frequency domain features (36~dimensions) via continuous wavelet transform.
To reduce dimensionality and mitigate overfitting, features are ranked by random forest feature importance~\cite{breiman2001} (200 trees, balanced class weights) and the top $k=10$ are retained for each training configuration.

\subsection{Domain Grouping}
\label{sec:domain}

Pairwise distances between all 87 subjects are computed on their mean feature vectors using three metrics:
\begin{itemize}
  \item \textbf{Maximum Mean Discrepancy (MMD)}~\cite{gretton2012}: measures distributional differences in a reproducing kernel Hilbert space.
  \item \textbf{Dynamic Time Warping (DTW)}~\cite{berndt1994}: aligns temporal sequences to capture pattern similarity.
  \item \textbf{Wasserstein Distance}~\cite{villani2009}: computes the minimum transport cost between distributions.
\end{itemize}
For each metric, a $K$-nearest-neighbor (KNN) score ($K=5$) is computed per subject, and subjects are split at the median into \textbf{in-domain} (44 subjects, most typical) and \textbf{out-domain} (43 subjects, most dissimilar).

\subsection{Experimental Design}
\label{sec:design}

The experiment follows a four-factor full-factorial design (Table~\ref{tab:design}).

\begin{table}[!t]
  \caption{Factorial Experiment Design}
  \label{tab:design}
  \centering
  \footnotesize
  \begin{tabular}{l c c p{2.9cm}}
    \toprule
    Factor & Sym. & Lvl. & Values \\
    \midrule
    Rebalancing & $R$ & 7 & BL, RUS/SMOTE/SW-SMOTE $r\!\in\!\{0.1,0.5\}$ \\
    Training mode & $M$ & 3 & Cross, Within, Mixed \\
    Distance & $D$ & 3 & MMD, DTW, Wasserstein \\
    Membership & $G$ & 2 & In-domain, Out-domain \\
    \bottomrule
  \end{tabular}
\end{table}

The $7\times 3\times 3\times 2 = 126$ factor combinations are each evaluated over 12 fixed random seeds, yielding $N = 1{,}512$ total observations.
Three primary evaluation metrics are used: \textbf{F2-score} (emphasizing recall for safety-critical detection), \textbf{AUROC} (overall discrimination), and \textbf{AUPRC} (precision--recall trade-off under class imbalance)~\cite{saito2015}.
Data are split per subject in chronological order: 60\% train, 20\% validation, 20\% test.

\textbf{Classifier and hyperparameter optimization.}
All 126 factor combinations use a random forest (RF) classifier~\cite{breiman2001} (scikit-learn \texttt{RandomForestClassifier}) whose hyperparameters are tuned independently per cell via Optuna~\cite{akiba2019} with 100 Tree-structured Parzen Estimator (TPE) trials and 3-fold stratified cross-validation, optimizing the F2-score.
Table~\ref{tab:hpo} summarizes the search space.
After Optuna selection, the classification threshold is swept over 1,001 equally spaced values in $[0,1]$, retaining the threshold that maximizes F2 on the validation set.
The final model is calibrated via Platt scaling~\cite{platt1999} (5-fold) on the combined train\,+\,validation data.

\begin{table}[!t]
  \caption{Random Forest Hyperparameter Search Space (Optuna, 100 TPE Trials)}
  \label{tab:hpo}
  \centering
  \footnotesize
  \begin{tabular}{l l c}
    \toprule
    Parameter & Range & Type \\
    \midrule
    \texttt{n\_estimators}      & 50--1{,}000             & int \\
    \texttt{max\_depth}         & \{None,10,20,30,50,100\} & cat \\
    \texttt{min\_samples\_split} & 2--100               & int \\
    \texttt{min\_samples\_leaf}  & 1--50                & int \\
    \texttt{max\_features}      & \{$\sqrt{p}$,$\log_2\!p$,0.1,\ldots,1.0\} & cat \\
    \texttt{max\_samples}       & \{None,0.5,0.7,0.9\} & cat \\
    \texttt{class\_weight}      & \{balanced,bal.\_sub.,None\} & cat \\
    \bottomrule
  \end{tabular}
\end{table}

\textbf{Computational cost.}
Each of the 1,512 training runs ($126\text{ cells}\times 12\text{ seeds}$) involves 100 Optuna trials with 3-fold CV, totalling $\approx\!4.5\times 10^5$ model fits.
Individual HPC jobs required 4 CPU cores, 8--10\,GB RAM, and 6--8 hours wall time; the entire experiment consumed $\approx\!10{,}000$ core-hours on the university cluster.

\textbf{Inference latency.}
A trained RF with $\leq\!1{,}000$ trees and $k=10$ input features produces a prediction in ${<}\,1$\,ms on a single CPU core, well within the 3-second window update interval.

\subsection{One-Factor-at-a-Time (OFAT) Analysis}
\label{sec:ofat}

Before the quantitative decomposition, we perform an OFAT analysis to build intuition.
For each target factor, all other factors are fixed at every possible combination, and the mean performance is plotted across the target factor's levels.
Each resulting line represents one fixed condition (averaged over 12 seeds), producing a ``spaghetti plot'' that reveals:
(i)~whether the factor's effect is consistent across conditions,
(ii)~how much the estimate varies depending on which conditions are fixed, and
(iii)~the magnitude of the factor's effect relative to the spread across conditions.

\subsection{Variance-Based Sensitivity Analysis (Sobol Indices)}
\label{sec:sobol}

To quantify the relative importance of each factor and their interactions, we adopt the Sobol--Hoeffding functional ANOVA decomposition~\cite{sobol1993,saltelli2008}.
For a model output $Y = f(X_1,\ldots,X_k)$, the total variance admits a unique orthogonal decomposition:
\begin{equation}
  V(Y) = \sum_{i=1}^{k} V_i + \sum_{i<j} V_{ij} + \cdots + V_{1,2,\ldots,k}
  \label{eq:anova}
\end{equation}
where $V_i = \mathrm{Var}_{X_i}\bigl[\mathbb{E}_{X_{\sim i}}(Y\mid X_i)\bigr]$ is the variance of the conditional expectation over factor $i$ alone, and $V_{ij}$ captures the joint effect of $(X_i, X_j)$ not explained by their individual contributions.
The first-order Sobol index normalises each contribution:
\begin{equation}
  S_i = \frac{V_i}{V(Y)}
  \label{eq:sobol1}
\end{equation}
and the total-order index aggregates the main effect with all interactions involving factor $i$:
\begin{equation}
  S_{Ti} = 1 - \frac{V_{\sim i}}{V(Y)}
  \label{eq:sobolT}
\end{equation}

The difference $S_{Ti} - S_i$ quantifies the fraction of variance attributable to interactions involving factor $i$.
Because our design is balanced, all sum-of-squares terms are computed exactly via marginal means without Monte Carlo sampling.
Confidence intervals (95\%) are obtained by percentile bootstrap ($B = 2{,}000$) resampling over the 12 seeds.

\subsection{Statistical Framework}
\label{sec:stats}

Due to non-normality (Shapiro--Wilk rejects normality in 45--71\% of cells), all analyses use non-parametric tests: Kruskal--Wallis $H$ for $k$-group comparisons~\cite{kruskal1952}, Friedman $\chi_F^2$ for paired rankings~\cite{friedman1937}, Mann--Whitney $U$ for pairwise unpaired tests, and Wilcoxon signed-rank for paired tests.
Effect sizes use Cliff's $\delta$~\cite{cliff1993} (${<}\,0.147$ negligible, ${<}\,0.33$ small, ${<}\,0.474$ medium, ${\geq}\,0.474$ large).
All test families are Bonferroni-corrected ($\alpha' = \alpha/m$, $\alpha = 0.05$).
For null claims (e.g., distance metric equivalence), we supplement with Bayes factors via BIC approximation~\cite{wagenmakers2007}, interpreted on the Jeffreys scale~\cite{jeffreys1961}.
A global permutation test ($B = 10{,}000$) validates the rebalancing effect.
Seed convergence analysis confirms ranking stability.

% =====================================================================
\section{Results}
\label{sec:results}

\subsection{Intuitive Factor Analysis via OFAT}
\label{sec:results-ofat}

Before the quantitative variance decomposition, we examine each factor's effect using OFAT spaghetti plots (Figs.~\ref{fig:ofat-R}--\ref{fig:ofat-G}).
Each thin line represents one fixed combination of the other three factors (averaged over 12 seeds); the bold black line shows the grand OFAT mean, and the gray band indicates $\pm 1$ SD across conditions.

\textbf{Rebalancing} ($R$; Fig.~\ref{fig:ofat-R}).
SMOTE and SW-SMOTE variants consistently outperform Baseline and RUS across all 18 fixed conditions ($D\times G\times M$).
The inter-condition spread ($\Delta = 0.08$--$0.55$ for F2-score) reflects strong mode-dependence: cross-domain conditions cluster at the bottom with small rebalancing gains, while within-domain and mixed conditions fan upward with large SMOTE-driven improvements.
This visual pattern already suggests a substantial $R\times M$ interaction.

\begin{figure}[!t]
  \centering
  \includegraphics[width=\columnwidth]{fig_s_ofat_condition.pdf}
  \caption{OFAT analysis for Rebalancing ($R$). Each line represents one of 18 fixed conditions ($D\times G\times M$), averaged over 12~seeds. Bold line: grand OFAT mean; gray band: $\pm 1$ SD. The large spread and mode-dependent fan confirm $R$'s strong, context-dependent effect.}
  \label{fig:ofat-R}
\end{figure}

\textbf{Distance} ($D$; Fig.~\ref{fig:ofat-D}).
The OFAT mean is flat across MMD, DTW, and Wasserstein, and individual condition lines remain nearly parallel.
This confirms that $D$ is negligible regardless of which other factors are fixed.

\begin{figure}[!t]
  \centering
  \includegraphics[width=\columnwidth]{fig_s_ofat_distance.pdf}
  \caption{OFAT analysis for Distance ($D$). Each line represents one of 42 fixed conditions ($R\times G\times M$). The flat trajectories confirm metric irrelevance.}
  \label{fig:ofat-D}
\end{figure}

\textbf{Mode} ($M$; Fig.~\ref{fig:ofat-M}).
Every condition shows a monotonic rise from cross-domain to within-domain/mixed.
At the Cross level, all lines converge near floor performance (F2\,$\approx 0.05$--$0.20$); at Within and Mixed, lines fan out according to rebalancing strategy---further evidence of the $R\times M$ interaction.

\begin{figure}[!t]
  \centering
  \includegraphics[width=\columnwidth]{fig_s_ofat_mode.pdf}
  \caption{OFAT analysis for Mode ($M$). The monotonic rise and fan-out pattern visualize the dominant main effect and $R\times M$ interaction.}
  \label{fig:ofat-M}
\end{figure}

\textbf{Membership} ($G$; Fig.~\ref{fig:ofat-G}).
The OFAT mean shows a slight upward slope from In-domain to Out-domain, but individual lines diverge in both directions.
The effect averages to near zero, confirming negligibility.

\begin{figure}[!t]
  \centering
  \includegraphics[width=\columnwidth]{fig_s_ofat_level.pdf}
  \caption{OFAT analysis for Membership ($G$). Divergent lines and near-flat mean confirm the negligible main effect.}
  \label{fig:ofat-G}
\end{figure}

\textbf{OFAT summary.}
The visual analysis establishes a clear hierarchy: $M$ and $R$ are the dominant factors with a visible interaction, while $D$ and $G$ are negligible.
The next section quantifies these observations.

\subsection{Quantitative Factor Decomposition via Sobol Indices}
\label{sec:results-sobol}

The Sobol--Hoeffding decomposition (Fig.~\ref{fig:sobol}) quantifies each factor's contribution to total variance.
The permutation test confirms a significant global rebalancing effect for all three metrics ($p < 0.001$).

\textbf{Mode} contributes the largest main effect ($S_M = 0.31$--$0.50$), followed by \textbf{Rebalancing} ($S_R = 0.24$--$0.29$).
Both factors participate in a substantial $R\times M$ interaction (21.2\% of F2-score variance, 13.8\% of AUROC variance).
The total-order indices show that Mode accounts for $S_{TM} = 0.48$--$0.66$ and Rebalancing for $S_{TR} = 0.40$--$0.46$ of total variance.

In contrast, \textbf{Distance} ($S_{TD} < 0.015$) and \textbf{Membership} ($S_{TG} < 0.031$) are negligible even when interactions are included.
Residual variance (seed-to-seed variation), denoted $S_{\varepsilon}$, accounts for 7.9\%--21.9\% of total variance.

Defining the systematic (non-residual) variance fraction, where $S_{\varepsilon}$ is the residual (seed-level) variance share:
\begin{equation}
  \frac{S_M + S_R + S_{R{\times}M}}{1 - S_{\varepsilon}}
  = \frac{0.823}{0.843}
  = 97.5\%\;\text{(F2)}
  \label{eq:systematic}
\end{equation}
where $S_M{=}0.368$, $S_R{=}0.243$, $S_{R{\times}M}{=}0.212$, and $S_{\varepsilon}{=}0.157$.
Analogous ratios are 95.8\% (AUROC) and 96.1\% (AUPRC), confirming that $M$, $R$, and their interaction account for ${>}\,95$\% of systematic variance across all three metrics.

\begin{figure}[!t]
  \centering
  \includegraphics[width=\columnwidth]{fig2_effect_hierarchy.pdf}
  \caption{Sobol sensitivity indices for the four factors. Solid bars: first-order indices ($S_i$); hatched bars: interaction contributions ($S_{Ti}-S_i$). Error bars: 95\% bootstrap CIs on $S_{Ti}$. Mode and Rebalancing dominate; Distance and Membership are negligible.}
  \label{fig:sobol}
\end{figure}

\textbf{$R\times M$ interaction concentration.}
The coupling between $R$ and $M$ is not only the largest interaction but effectively the \emph{only} one.
The interaction concentration ratio $C_{R\times M}^{(i)} = S_{R\times M}/(S_{Ti}-S_i)$---the share of factor $i$'s total interaction attributable to the $R\times M$ term---yields:

\begin{table}[!t]
  \caption{Interaction Concentration Ratios}
  \label{tab:concentration}
  \centering
  \begin{tabular}{l c c}
    \toprule
    Metric & $C_{R\times M}^{(M)}$ & $C_{R\times M}^{(R)}$ \\
    \midrule
    F2-score & 94\% & 96\% \\
    AUROC    & 87\% & 88\% \\
    AUPRC    & 88\% & 89\% \\
    \bottomrule
  \end{tabular}
\end{table}

The factorial design's $2^k - k - 1 = 11$ interaction terms collapse into a single dominant pairwise effect.

\textbf{Distance metric equivalence---Bayesian confirmation.}
Because frequentist tests can only fail to reject $H_0$, we supplement with Bayes factors:

\begin{table}[!t]
  \caption{Bayes Factors for Distance Metric Equivalence}
  \label{tab:bayes}
  \centering
  \begin{tabular}{l c c}
    \toprule
    Metric & $BF_{01}$ & Interpretation \\
    \midrule
    F2-score & 767 & Extreme evidence for $H_0$ \\
    AUROC    &  71 & Very strong evidence for $H_0$ \\
    AUPRC    & 381 & Extreme evidence for $H_0$ \\
    \bottomrule
  \end{tabular}
\end{table}

All Bayes factors provide ``very strong'' to ``extreme'' evidence on the Jeffreys scale~\cite{jeffreys1961} that the distance metric has no effect, affirming the Sobol result ($S_{TD} < 0.015$).

\begin{figure}[!t]
  \centering
  \includegraphics[width=\columnwidth]{fig5_distance_violin.pdf}
  \caption{Performance distributions by distance metric (within-domain and mixed modes). The three metrics produce virtually indistinguishable distributions ($|\delta|<0.15$ for all pairwise comparisons).}
  \label{fig:violin}
\end{figure}

\subsection{Inter-Factor Interactions}
\label{sec:results-interactions}

The Sobol decomposition identified a single dominant interaction ($R\times M$).
This section characterizes its structure using three complementary analyses.

\subsubsection{Performance by Strategy and Mode}

Table~\ref{tab:pooled} shows mean performance by rebalancing strategy (all modes pooled), and Table~\ref{tab:permode} disaggregates by training mode.
Within-domain and Mixed produce statistically equivalent results ($\delta < 0.05$), while Cross-domain collapses to near-chance levels ($\delta > 0.83$ vs.\ Within).

\begin{table}[!t]
  \caption{Mean $\pm$ SD Performance by Rebalancing Strategy (All Modes Pooled)}
  \label{tab:pooled}
  \centering
  \footnotesize
  \begin{tabular}{l c c c}
    \toprule
    Strategy & F2 & AUROC & AUPRC \\
    \midrule
    Baseline        & $.215{\pm}.057$ & $.631{\pm}.123$ & $.135{\pm}.168$ \\
    RUS $r\!=\!.1$  & $.184{\pm}.049$ & $.594{\pm}.094$ & $.108{\pm}.122$ \\
    RUS $r\!=\!.5$  & $.162{\pm}.036$ & $.564{\pm}.064$ & $.072{\pm}.052$ \\
    SMOTE $r\!=\!.1$  & $.346{\pm}.158$ & $.765{\pm}.176$ & $.407{\pm}.286$ \\
    SMOTE $r\!=\!.5$  & $.376{\pm}.205$ & $.755{\pm}.172$ & $.394{\pm}.276$ \\
    SW-SM. $r\!=\!.1$  & $\mathbf{.416}{\pm}.243$ & $\mathbf{.765}{\pm}.183$ & $\mathbf{.447}{\pm}.337$ \\
    SW-SM. $r\!=\!.5$  & $.304{\pm}.233$ & $.745{\pm}.166$ & $.274{\pm}.223$ \\
    \bottomrule
  \end{tabular}
\end{table}

\begin{table*}[!t]
  \caption{Mean Performance by Rebalancing Strategy, Disaggregated by Training Mode}
  \label{tab:permode}
  \centering
  \begin{tabular}{l c c c c c}
    \toprule
    Strategy & F2 (Cross) & F2 (Within) & F2 (Mixed) & AUROC (Within) & AUPRC (Within) \\
    \midrule
    Baseline              & 0.160 & 0.215 & 0.268 & 0.633 & 0.116 \\
    RUS $r\!=\!0.1$       & 0.165 & 0.178 & 0.208 & 0.625 & 0.122 \\
    RUS $r\!=\!0.5$       & 0.156 & 0.159 & 0.170 & 0.603 & 0.097 \\
    SMOTE $r\!=\!0.1$     & 0.138 & 0.448 & 0.452 & 0.898 & 0.600 \\
    SMOTE $r\!=\!0.5$     & 0.114 & 0.499 & 0.514 & 0.882 & 0.562 \\
    SW-SMOTE $r\!=\!0.1$  & 0.101 & $\mathbf{0.558}$ & $\mathbf{0.587}$ & $\mathbf{0.903}$ & $\mathbf{0.648}$ \\
    SW-SMOTE $r\!=\!0.5$  & 0.042 & 0.468 & 0.402 & 0.862 & 0.386 \\
    \bottomrule
  \end{tabular}
\end{table*}

\begin{figure*}[!t]
  \centering
  \includegraphics[width=\textwidth]{fig4_strategy_comparison.pdf}
  \caption{Performance distributions of the 7 strategies by training mode (rows: F2, AUROC, AUPRC; columns: Cross, Within, Mixed). In Cross-domain, all strategies compress to a narrow low-performance band. In Within/Mixed, SMOTE-based strategies separate sharply from Baseline and RUS.}
  \label{fig:strategy}
\end{figure*}

\subsubsection{Strategy Ranking Reversal}

The $R\times M$ interaction manifests as a qualitative reversal: the \emph{identity} of the best strategy changes across modes.
Per-mode Friedman tests (all $p < 10^{-5}$) show RUS $r\!=\!0.1$ ranks first in cross-domain, while SW-SMOTE $r\!=\!0.1$ ranks first in within-domain and mixed.
Cross-mode Spearman $\rho$ confirms the reversal:
\begin{equation}
  \rho_{\mathrm{cross,within}} = -0.74\;\text{to}\;-0.89
  \label{eq:reversal}
\end{equation}
while within- vs.\ mixed yields $\rho \geq +0.99$ ($p < 0.001$).
Kendall's $W$ across the three modes is $0.17$--$0.22$ (non-significant), confirming that modes do \emph{not} agree on strategy rankings.

\begin{figure}[!t]
  \centering
  \includegraphics[width=\columnwidth]{fig_ranking_reversal.pdf}
  \caption{Strategy ranking reversal---bump chart of Friedman mean ranks across training modes. In Cross-domain, RUS variants rank highest; in Within/Mixed, SMOTE-based strategies dominate. The near-perfect overlap of Within and Mixed ($\rho \geq 0.99$) contrasts with the complete inversion of Cross vs.\ Within ($\rho = -0.74$ to $-0.89$).}
  \label{fig:bump}
\end{figure}

The practical consequence is that a practitioner who selects a rebalancing strategy based on cross-domain results would deploy the \emph{worst} strategy for within-domain operation, and vice versa (Table~\ref{tab:permode}).

\subsubsection{Domain Gap Reversal}

An unexpected finding is that the domain gap reverses in within-domain and mixed training: out-domain subjects sometimes outperform in-domain subjects ($\Delta > 0$).
Fig.~\ref{fig:domshift} visualizes this through diverging bars for each $R\times M\times G$ cell.
Blue bars (positive $\Delta$) indicate out-domain outperformance.

\begin{figure}[!t]
  \centering
  \includegraphics[width=\columnwidth]{fig7_domain_shift_reversal.pdf}
  \caption{Domain gap direction ($\Delta = \mathrm{out} - \mathrm{in}$) by $R\times M$. Blue bars = out-domain outperforms in-domain (gap reversal). The prevalence of blue bars in Mixed mode confirms that domain membership does not cause systematic degradation.}
  \label{fig:domshift}
\end{figure}

The Sobol decomposition confirms that while the direction reversal is qualitatively notable, its magnitude is small: $S_{G\times M} = 0.006$--$0.009$ (0.6--0.9\% of total variance), an order of magnitude below $S_{R\times M}$.

\subsection{Robustness Validation}
\label{sec:robustness}

\textbf{Cross-metric concordance.}
Kendall's $W = 0.643$ across scores; AUROC--AUPRC $\rho = 0.929$.
Results are consistent across~5 metrics and~2 sampling ratios (91\% directional agreement).

\textbf{Post-hoc power.}
The pooled design ($n = 36$) detects $|\delta| \geq 0.53$, well below all reported large effects ($|\delta| > 0.8$).

\textbf{Strategy clusters.}
Nemenyi post-hoc tests~\cite{nemenyi1963} (CD\,=\,2.600) identify two statistically separated clusters: oversampling methods (mean ranks 2.76--3.83) vs.\ Baseline/RUS (ranks 4.12--5.64).

\textbf{Seed convergence.}
Ranking stability ($\sigma_{\mathrm{rank}}$) decreases monotonically with the number of seeds, reaching 0 (F2, AUPRC) or 0.147 (AUROC) by $k = 11$ of 12 seeds (Fig.~\ref{fig:seed}), confirming sufficient seed count.

\begin{figure}[!t]
  \centering
  \includegraphics[width=\columnwidth]{fig8_seed_convergence.pdf}
  \caption{Ranking stability ($\sigma_{\mathrm{rank}}$) as a function of seed subset size $k$. Monotonic convergence confirms that $n = 12$ seeds is sufficient.}
  \label{fig:seed}
\end{figure}

\subsection{Classifier-Independence Validation (XGBoost)}
\label{sec:results-xgboost}

To assess whether the variance decomposition reflects data-level structure rather than classifier-specific artifacts, we replicate the full 1,512-cell factorial experiment with XGBoost~\cite{chen2016}, an architecturally distinct gradient boosting classifier.
Whereas random forest reduces variance via bagging and random feature subsets, XGBoost reduces bias through sequential residual correction---a fundamentally different inductive bias.
Each XGBoost configuration uses Optuna~\cite{akiba2019} with 100 trials for hyperparameter optimization, identical to the RF protocol.

Table~\ref{tab:sobol-xgb} compares the Sobol indices.
The factor hierarchy is preserved: Mode remains the dominant main effect ($S_M^{\mathrm{XGB}} = 0.50$ vs.~$S_M^{\mathrm{RF}} = 0.37$), Rebalancing is second ($S_R^{\mathrm{XGB}} = 0.24$ vs.~$S_R^{\mathrm{RF}} = 0.24$), and the $R\times M$ interaction is substantial ($S_{R\times M}^{\mathrm{XGB}} = 0.11$ vs.~$S_{R\times M}^{\mathrm{RF}} = 0.21$).
Distance and Membership remain negligible ($S_D^{\mathrm{XGB}} = 0.002$, $S_G^{\mathrm{XGB}} = 0.010$).
The systematic concentration ratio $\frac{S_M + S_R + S_{R\times M}}{1 - S_{\varepsilon}} = 98.6\%$ (XGB) remains above 90\%, confirming that the dominance of $M$, $R$, and $R\times M$ is not an artifact of the RF learning algorithm.

\begin{table}[!t]
  \caption{Sobol First-Order Indices: RF vs.\ XGBoost (F2-Score)}
  \label{tab:sobol-xgb}
  \centering
  \begin{tabular}{l c c}
    \toprule
    Factor & RF $S_i$ & XGBoost $S_i$ \\
    \midrule
    Mode ($M$)          & 0.368 & 0.501 \\
    Rebalancing ($R$)   & 0.243 & 0.235 \\
    $R \times M$        & 0.212 & 0.113 \\
    Distance ($D$)      & 0.005 & 0.002 \\
    Membership ($G$)    & 0.010 & 0.010 \\
    Residual ($\varepsilon$) & 0.157 & 0.139 \\
    \bottomrule
  \end{tabular}
\end{table}

The ranking reversal pattern also replicates.
In within-domain and mixed modes, SW-SMOTE $r\!=\!0.1$ achieves the best F2 (0.758 within, 0.781 mixed), while in cross-domain mode all strategies collapse to F2\,$= 0.18$--$0.22$.
The Spearman rank correlation between within-domain and mixed mode strategy rankings is $\rho = 0.96$ ($p < 0.001$), confirming the same Within--Mixed equivalence observed with RF.
Fig.~\ref{fig:sobol-comparison} visualizes the side-by-side comparison.

\begin{figure}[!t]
  \centering
  \includegraphics[width=\columnwidth]{fig9_sobol_comparison.pdf}
  \caption{Sobol index comparison between RF (blue) and XGBoost (orange). Solid bars: first-order indices; hatched: interaction contributions. The factor hierarchy ($M > R \gg D \approx G$) is preserved across both classifier families.}
  \label{fig:sobol-comparison}
\end{figure}

% =====================================================================
\section{Discussion}
\label{sec:discussion}

\subsection{Why Training Mode and Rebalancing Dominate}
\label{sec:disc-dominate}

The Sobol decomposition reveals a remarkably concentrated variance structure: $M$, $R$, and $R\times M$ jointly explain ${>}\,95$\% of systematic variance (Eq.~\ref{eq:systematic}), while $D$ and $G$ together contribute ${<}\,5$\%.
This concentration has a data-geometric explanation.

In \textbf{within-domain and mixed modes}, each subject's own data is present in the training set, providing a subject-specific decision boundary.
The performance bottleneck is class imbalance: minority-class (drowsy) epochs are scarce, and SMOTE generates synthetic examples along the minority-class manifold, enriching the boundary without discarding majority-class information:
\begin{equation}
  \Delta\mathrm{F2} = 0.558 - 0.215 = {+}0.343\;({+}159\%)
  \label{eq:delta-f2}
\end{equation}

RUS, by contrast, discards majority-class samples ($\delta_{\mathrm{RUS\,vs\,BL}} = -0.84$), destroying the information that calibrates subject-specific patterns.

In \textbf{cross-domain mode}, target-subject data is entirely absent.
The bottleneck shifts from class imbalance to domain mismatch: oversampling amplifies source-domain patterns that may not transfer, while RUS's information loss is less damaging because the majority-class signal was already misaligned.
All strategies collapse to near-chance levels (F2\,$= 0.04$--$0.17$).
This mode-dependent bottleneck shift is the mechanistic origin of the $R\times M$ interaction and the ranking reversal (Eq.~\ref{eq:reversal}).

For comparison, switching distance metric yields $|\Delta\mathrm{F2}| < 0.01$---two orders of magnitude below the rebalancing effect, consistent with $S_{TR}/S_{TD} > 27\times$.

\subsection{Why Distance Metrics Are Irrelevant}
\label{sec:disc-distance}

The finding that three fundamentally different metrics---MMD (distributional), DTW (temporal), and Wasserstein (geometric)---produce equivalent downstream performance ($S_{TD} < 0.015$, $BF_{01} = 71$--$767$) is particularly notable because they generate genuinely different subject rankings under normalised features:

\begin{table}[!t]
  \caption{Pairwise Spearman Correlations Between Distance Metrics}
  \label{tab:spearman}
  \centering
  \begin{tabular}{l c}
    \toprule
    Pair & Spearman $\rho$ \\
    \midrule
    MMD vs.\ Wasserstein & 0.757 \\
    MMD vs.\ DTW         & 0.116 (n.s.) \\
    Wasserstein vs.\ DTW & 0.444 \\
    \bottomrule
  \end{tabular}
\end{table}

DTW captures distributional properties genuinely distinct from MMD and Wasserstein ($\rho = 0.12$, $p = 0.29$), and 42.5\% of subjects switch domain groups across metrics.
Yet classification performance is unchanged.
Three mechanisms explain this decoupling:

\begin{enumerate}
  \item \textbf{Physical coupling reduces effective dimensionality.}
  The bicycle model (Eq.~\ref{eq:bicycle}) couples $\delta$, $\dot{\delta}$, $a_y$, and $e_{\mathrm{lane}}$ through a causal chain, with only $a_x$ independent.
  PCA confirms a 3:1 compression (45 components explain 95\% of variance from 135 features).

  \item \textbf{Weak inter-subject discrimination.}
  The intraclass correlation coefficient computed on our 87-subject dataset following the ICC(1,1) formulation of Shrout and Fleiss~\cite{shrout1979} yields ICC\,=\,0.111, meaning 88.9\% of feature variance is intra-subject.
  When subjects' positions in feature space are this ``blurred,'' the coarse binary partition at the median absorbs ranking differences without substantially altering the downstream training set.

  \item \textbf{Rebalancing absorption.}
  The Sobol ratio $S_{TR}/S_{TD} > 27\times$ (Eq.~\ref{eq:systematic}) confirms that rebalancing shifts the decision boundary so dramatically that grouping differences are overwhelmed:
  \begin{equation}
    \frac{S_{TR}}{S_{TD}} \approx \frac{0.40\text{--}0.46}{{<}\,0.015} = 27\text{--}31\times
    \label{eq:ratio}
  \end{equation}
\end{enumerate}

\subsection{Practical Implications}
\label{sec:disc-practical}

The findings prescribe a clear decision hierarchy for practitioners deploying vehicle-dynamics-based DDD systems:

\begin{enumerate}
  \item \textbf{Training mode (highest priority):}
  Use within-domain or mixed training whenever target-subject data is available.
  This is the single largest factor ($S_M = 0.31$--$0.50$).

  \item \textbf{Class rebalancing (second priority):}
  Apply SMOTE-based oversampling at conservative ratio ($r = 0.1$).
  SW-SMOTE $r\!=\!0.1$ achieves the best overall performance (+159\% F2, +457\% AUPRC over baseline).
  If only cross-domain training is feasible, no rebalancing strategy provides substantial improvement.

  \item \textbf{Distance metric and domain membership (negligible):}
  Use any convenient metric for domain grouping.
  The choice has no measurable impact on classification performance ($BF_{01} > 70$).

  \item \textbf{Real-time feasibility:}
  RF inference with 10 features and $\leq\!1{,}000$ trees requires ${<}\,1$\,ms per window on a single CPU, comfortably meeting the 3-second sliding window update rate without specialized hardware.

  \item \textbf{Bridging the cross-domain gap.}
  Our results show that cross-domain performance remains near chance for all rebalancing strategies (F2\,$= 0.04$--$0.17$), highlighting the need for complementary techniques.
  Three promising directions exist.
  First, unsupervised domain adaptation methods~\cite{pan2010,weiss2016} that align source and target feature distributions without labeled target data could raise the cross-domain floor without requiring per-driver labels.
  Second, few-shot or zero-shot personalization~\cite{kim2025}, in which a small number of unlabeled target-driver epochs are used for distributional calibration, could approximate within-domain performance at a fraction of the data cost.
  Third, continual learning strategies that incrementally update a deployed model as new driver data accumulates would provide a practical path from cross-domain to within-domain operation over time.
  Integrating such approaches within the factorial framework presented here would allow practitioners to quantify their added value relative to the dominant $M$ and $R$ factors and to identify which combination minimizes the deployment data requirement.
\end{enumerate}

\subsection{Comparison with Prior Work}
\label{sec:disc-comparison}

Our finding that class imbalance handling is the dominant design factor aligns with Taylor et~al.~\cite{taylor2021}, who showed that data-related hyperparameters dominate architectural choices in deep learning.
However, our study is the first to demonstrate this in the DDD context and, crucially, to reveal the interaction with training mode that causes strategy ranking reversal.
Prior DDD studies that compare rebalancing methods under a single training mode~\cite{he2024} risk selecting a suboptimal strategy for deployment.

The best within-domain performance (F2\,=\,0.558, AUROC\,=\,0.903) is competitive with recent vehicle-dynamics-based methods.
Table~\ref{tab:comparison} positions this work relative to prior DDD studies.
Direct performance comparison is limited because prior studies use different datasets, modalities, and evaluation protocols; the primary contribution here is the factorial decomposition framework, not absolute classification performance.
Nevertheless, the within-domain AUROC (0.903) is competitive with Wang et~al.~\cite{wang2022} (0.91 on the same dataset using an LSTM with a single fixed split), confirming that the RF pipeline optimized via Optuna provides a credible experimental baseline.

\begin{table*}[!t]
  \caption{Comparison with Prior DDD Studies. ``Factors'' Indicates the Number of Design Factors Systematically Varied; ``Interaction'' Indicates Whether Inter-Factor Interactions Were Quantified. Studies Using the Same 87-Subject Dataset~\cite{aygun2024} Are Marked with~$\dagger$.}
  \label{tab:comparison}
  \centering
  \footnotesize
  \begin{tabular}{l c l l c c c c}
    \toprule
    Study & Year & Classif. & Modal. & Subj. & Fact. & Inter. & AUROC \\
    \midrule
    Arefnezhad et~al.~\cite{arefnezhad2019}   & 2019 & SVM/NN    & Vehicle & 20 & 0            & No               & 0.87 \\
    Wang et~al.~\cite{wang2022}$^{\dagger}$    & 2022 & LSTM      & Vehicle & 87 & 0            & No               & 0.91 \\
    Zhao et~al.~\cite{zhao2012}                & 2012 & SVM       & Vehicle & 15 & 0            & No               & 0.82 \\
    He et~al.~\cite{he2024}                    & 2024 & GA-CNN    & EEG     & 12 & 1 (imbalance)& No               & 0.89 \\
    Aravinth et~al.~\cite{aravinth2025}        & 2025 & Transfer  & EEG     & 20 & 1 (domain)   & No               & 0.88 \\
    Kim et~al.~\cite{kim2025}                  & 2025 & Mixup     & EEG     & 27 & 1 (domain)   & No               & 0.92 \\
    Luo et~al.~\cite{luo2024}                  & 2024 & DA-Net    & EEG     & 15 & 1 (domain)   & No               & 0.85 \\
    Li et~al.~\cite{li2026}                    & 2026 & MADNet    & EEG     & 24 & 1 (domain)   & No               & 0.90 \\
    Fonseca \& Ferreira~\cite{fonseca2025}     & 2025 & Review    & Multi   & --- & ---          & ---              & --- \\
    \textbf{Ours (RF)}$^{\dagger}$              & \textbf{2026} & \textbf{RF} & \textbf{Vehicle} & \textbf{87} & \textbf{4} ($R$,$M$,$D$,$G$) & $R{\times}M$ \textbf{Sobol} & \textbf{0.903} \\
    \textbf{Ours (XGB)}$^{\dagger}$             & \textbf{2026} & \textbf{XGBoost} & \textbf{Vehicle} & \textbf{87} & \textbf{4} ($R$,$M$,$D$,$G$) & $R{\times}M$ \textbf{Sobol} & \textbf{0.877} \\
    \bottomrule
  \end{tabular}
\end{table*}

The distance metric irrelevance contrasts with the implicit assumption in domain adaptation literature~\cite{pan2010,kim2025} that domain definition substantially impacts downstream performance.
Our vehicle dynamics explanation---particularly the ICC\,=\,0.111~\cite{shrout1979} finding---suggests that for vehicle-dynamics-based DDD, the inter-subject signal is inherently too weak to support meaningful domain partitioning, regardless of the metric used.

\subsection{Limitations}
\label{sec:limitations}

\begin{enumerate}
  \item \textbf{Two classifier families tested.}
  The factorial experiment was conducted with both RF~\cite{breiman2001} and XGBoost~\cite{chen2016}, representing bagging and boosting paradigms (\S\ref{sec:results-xgboost}).
  The factor hierarchy $S_{TM} > S_{TR} \gg S_{TD} \approx S_{TG}$ is consistent across both, supporting classifier independence.
  However, extension to fundamentally different architectures (SVM, LSTM, Transformer) remains an open question.

  \item \textbf{Deterministic data split.}
  The train/test partition is deterministic; random seeds vary only model initialisation and resampling.

  \item \textbf{Simulator data.}
  All data come from a laboratory driving simulator.
  Crucially, the five vehicle-dynamics signals used here ($\delta$, $\dot{\delta}$, $a_y$, $a_x$, $e_{\mathrm{lane}}$) are governed by the same bicycle model (Eqs.~\ref{eq:bicycle}--\ref{eq:lane}) in both simulated and real vehicles; the physical coupling that produces feature redundancy and low ICC (0.111) is intrinsic to vehicle dynamics, not an artifact of the simulator.
  Nevertheless, on-road conditions introduce additional variability (road surface, weather, traffic interactions) that may alter the absolute variance shares.
  Validation on naturalistic driving datasets---such as the SHRP~2 corpus---is an essential next step to confirm that the factor hierarchy $M > R \gg D \approx G$ transfers to real-world conditions.

  \item \textbf{Sample size per cell.}
  With $n = 12$, only large effects ($|\delta| > 0.53$) are detectable after Bonferroni correction.
  Subtle distance metric effects may exist below the detection threshold, though the Bayes factors ($BF_{01} > 70$) argue strongly against this possibility.

  \item \textbf{Feature selection scope.}
  Only the top $k = 10$ features (by RF importance~\cite{breiman2001}) are used per configuration; XGBoost's built-in gain-based importance produces a largely overlapping top-10 set, suggesting that the feature selection is not classifier-specific.
  While this aggressive reduction mitigates overfitting for the per-cell sample sizes in this design, a sensitivity analysis over $k \in \{5, 10, 20, 50\}$ would clarify whether the Sobol hierarchy is robust to feature set size.

  \item \textbf{Binary labeling threshold.}
  KSS\,$\leq 5$ vs.\ KSS\,$\geq 8$ excludes intermediate drowsiness (KSS\,6--7), which may limit applicability to graded drowsiness detection scenarios.
\end{enumerate}

% =====================================================================
\section{Conclusion}
\label{sec:conclusion}

This study presents the first factorial experimental analysis of drowsy driving detection system design, decomposing performance variance across four factors using Sobol--Hoeffding sensitivity analysis over 3,024 observations (1,512 RF + 1,512 XGBoost) from 87 drivers.
The key findings are:

\begin{enumerate}
  \item \textbf{Training mode and rebalancing dominate.}
  $M$ ($S_{TM} = 0.48$--$0.66$) and $R$ ($S_{TR} = 0.40$--$0.46$), including their interaction ($S_{R\times M} = 0.12$--$0.21$), account for ${>}\,95$\% of systematic variance.
  SW-SMOTE $r\!=\!0.1$ achieves the best within-domain performance (F2\,=\,0.558, AUROC\,=\,0.903 for RF; F2\,=\,0.758, AUROC\,=\,0.877 for XGBoost).

  \item \textbf{A strong interaction causes strategy ranking reversal.}
  The optimal rebalancing strategy depends on training mode: RUS is effective only in cross-domain settings, while SMOTE dominates in within-domain and mixed modes ($\rho_{\mathrm{cross,within}} = -0.74$ to $-0.89$).
  This reversal pattern is consistent across both RF and XGBoost.

  \item \textbf{Distance metric and domain membership are irrelevant.}
  $D$ ($S_{TD} < 0.015$; $BF_{01} = 71$--$767$) and $G$ ($S_{TG} < 0.031$) contribute negligibly.
  Vehicle dynamics coupling, weak inter-subject discrimination (ICC\,=\,0.111), and rebalancing absorption ($S_{TR}/S_{TD} > 27\times$) explain this irrelevance.

  \item \textbf{Classifier independence.}
  A full replication with XGBoost~\cite{chen2016}---an architecturally distinct gradient boosting classifier---preserves the factor hierarchy $S_{TM} > S_{TR} \gg S_{TD} \approx S_{TG}$, confirming that the variance structure reflects inherent data properties rather than classifier-specific artifacts.
\end{enumerate}

For practitioners, these results prescribe a simple decision rule: prioritize within-domain training and SMOTE-based class rebalancing; distance metric and domain grouping strategy require no optimization.
The factorial sensitivity analysis framework introduced here can be extended to other safety-critical detection tasks where multiple design factors interact.

\section*{Data and Code Availability}
The driving simulator dataset is publicly available at Harvard Dataverse (DOI: \href{https://doi.org/10.7910/DVN/HMZ5RG}{10.7910/DVN/HMZ5RG})~\cite{aygun2024}.
Analysis code and trained models will be released upon publication.

% =====================================================================
% REFERENCES
\bibliographystyle{IEEEtran}
% Inline bib for self-contained compilation
\begin{thebibliography}{43}

\bibitem{who2018}
World Health Organization, ``Global Status Report on Road Safety,'' 2018.

\bibitem{chawla2002}
N.~V. Chawla, K.~W. Bowyer, L.~O. Hall, and W.~P. Kegelmeyer, ``SMOTE: Synthetic minority over-sampling technique,'' \emph{J. Artif. Intell. Res.}, vol.~16, pp.~321--357, 2002.

\bibitem{he2009}
H.~He and E.~A. Garcia, ``Learning from imbalanced data,'' \emph{IEEE Trans. Knowl. Data Eng.}, vol.~21, no.~9, pp.~1263--1284, 2009.

\bibitem{pan2010}
S.~J. Pan and Q.~Yang, ``A survey on transfer learning,'' \emph{IEEE Trans. Knowl. Data Eng.}, vol.~22, no.~10, pp.~1345--1359, 2010.

\bibitem{weiss2016}
K.~Weiss, T.~M. Khoshgoftaar, and D.~D. Wang, ``A survey of transfer learning,'' \emph{J. Big Data}, vol.~3, no.~1, pp.~1--40, 2016.

\bibitem{arefnezhad2019}
S.~Arefnezhad \emph{et~al.}, ``Driver drowsiness estimation using EEG signals with a dynamical encoder--decoder,'' \emph{IET Intell. Transp. Syst.}, vol.~13, no.~2, pp.~301--310, 2019.

\bibitem{wang2022}
X.~Wang \emph{et~al.}, ``Real-time detection of driver drowsiness using LSTM,'' \emph{Sensors}, vol.~22, no.~13, p.~4904, 2022.

\bibitem{zhao2012}
C.~Zhao \emph{et~al.}, ``Driver drowsiness detection using continuous wavelet transform,'' in \emph{Proc. IEEE Int. Conf. Signal Process.}, 2012, pp.~1--4.

\bibitem{fonseca2025}
T.~Fonseca and S.~Ferreira, ``Drowsiness detection in drivers: A systematic review of deep learning-based models,'' \emph{Appl. Sci.}, vol.~15, no.~16, p.~9018, 2025.

\bibitem{alquraishi2024}
M.~S. Al-Quraishi \emph{et~al.}, ``Technologies for detecting and monitoring drivers' states: A systematic review,'' \emph{Heliyon}, vol.~10, no.~21, e38592, 2024.

\bibitem{he2024}
L.~He, L.~Zhang, Q.~Sun, and X.~T. Lin, ``A generative adaptive convolutional neural network with attention mechanism for driver fatigue detection with class-imbalanced and insufficient data,'' \emph{Behav. Brain Res.}, vol.~462, p.~114878, 2024.

\bibitem{aravinth2025}
S.~S. Aravinth, G.~M. Nagamani, C.~K. Kumar, and A.~Lasisi, ``Dynamic cross-domain transfer learning for driver fatigue monitoring,'' \emph{Sci. Rep.}, vol.~15, p.~92701, 2025.

\bibitem{kim2025}
D.~Y. Kim, D.~K. Han, and J.~H. Jeong, ``Calibration-free driver drowsiness classification with prototype-based multi-domain mixup,'' \emph{IEEE Trans. Intell. Transp. Syst.}, 2025.

\bibitem{luo2024}
Y.~Luo, W.~Liu, H.~Li, Y.~Lu, and B.~L. Lu, ``A cross-scenario and cross-subject domain adaptation method for driving fatigue detection,'' \emph{J. Neural Eng.}, vol.~21, no.~4, p.~046019, 2024.

\bibitem{li2026}
S.~Li, H.~Wang, X.~Liu, and R.~Zhu, ``MADNet: EEG modality augmentation for domain-generalized cross-subject drowsiness diagnosis,'' \emph{ACM Trans. Multimedia Comput. Commun. Appl.}, 2026.

\bibitem{sobol1993}
I.~M. Sobol', ``Sensitivity estimates for nonlinear mathematical models,'' \emph{Math. Model. Comput. Exp.}, vol.~1, no.~4, pp.~407--414, 1993.

\bibitem{tosin2020}
M.~Tosin, A.~M.~A. C\^{o}rtes, and A.~Cunha~Jr., ``A tutorial on Sobol' global sensitivity analysis applied to biological models,'' in \emph{Networks in Systems Biology}, Springer, 2020, pp.~93--118.

\bibitem{taylor2021}
R.~Taylor, V.~Ojha, I.~Martino, and G.~Nicosia, ``Sensitivity analysis for deep learning: Ranking hyper-parameter influence,'' in \emph{Proc. IEEE Symp. Ser. Comput. Intell.}, 2021, pp.~1--10.

\bibitem{kuhnt2022}
S.~Kuhnt and A.~Kalka, ``Global sensitivity analysis for the interpretation of machine learning algorithms,'' in \emph{Artificial Intelligence, Big Data and Data Science in Statistics}, Springer, 2022, pp.~135--164.

\bibitem{rajamani2012}
R.~Rajamani, \emph{Vehicle Dynamics and Control}, 2nd~ed. New York, NY, USA: Springer, 2012.

\bibitem{atiquzzaman2018}
M.~Atiquzzaman \emph{et~al.}, ``Real-time detection of drivers' texting and eating behavior based on vehicle dynamics,'' \emph{Transp. Res. Part~F}, vol.~58, pp.~594--604, 2018.

\bibitem{gretton2012}
A.~Gretton \emph{et~al.}, ``A kernel two-sample test,'' \emph{J. Mach. Learn. Res.}, vol.~13, no.~1, pp.~723--773, 2012.

\bibitem{berndt1994}
D.~J. Berndt and J.~Clifford, ``Using dynamic time warping to find patterns in time series,'' in \emph{AAAI Workshop Knowl. Discov. Databases}, 1994, pp.~359--370.

\bibitem{villani2009}
C.~Villani, \emph{Optimal Transport: Old and New}. Berlin, Germany: Springer, 2009.

\bibitem{saito2015}
T.~Saito and M.~Rehmsmeier, ``The precision--recall plot is more informative than the ROC plot when evaluating binary classifiers on imbalanced datasets,'' \emph{PLOS ONE}, vol.~10, no.~3, e0118432, 2015.

\bibitem{saltelli2008}
A.~Saltelli, M.~Ratto, T.~Andres, F.~Campolongo, J.~Cariboni, D.~Gatelli, M.~Saisana, and S.~Tarantola, \emph{Global Sensitivity Analysis: The Primer}. Chichester, UK: Wiley, 2008.

\bibitem{cliff1993}
N.~Cliff, ``Dominance statistics: Ordinal analyses to answer ordinal questions,'' \emph{Psychol. Bull.}, vol.~114, no.~3, pp.~494--509, 1993.

\bibitem{wagenmakers2007}
E.-J. Wagenmakers, ``A practical solution to the pervasive problems of $p$ values,'' \emph{Psychon. Bull. Rev.}, vol.~14, no.~5, pp.~779--804, 2007.

\bibitem{jeffreys1961}
H.~Jeffreys, \emph{Theory of Probability}, 3rd~ed. Oxford, UK: Oxford Univ. Press, 1961.

\bibitem{kruskal1952}
W.~H. Kruskal and W.~A. Wallis, ``Use of ranks in one-criterion variance analysis,'' \emph{J. Amer. Stat. Assoc.}, vol.~47, no.~260, pp.~583--621, 1952.

\bibitem{friedman1937}
M.~Friedman, ``The use of ranks to avoid the assumption of normality,'' \emph{J. Amer. Stat. Assoc.}, vol.~32, no.~200, pp.~675--701, 1937.

\bibitem{nemenyi1963}
P.~Nemenyi, ``Distribution-free multiple comparisons,'' Ph.D. dissertation, Princeton Univ., Princeton, NJ, USA, 1963.

\bibitem{masson2011}
M.~E.~J. Masson, ``A tutorial on a practical Bayesian alternative to null-hypothesis significance testing,'' \emph{Behav. Res. Methods}, vol.~43, no.~3, pp.~679--690, 2011.

\bibitem{aygun2024}
A.~Aygun \emph{et~al.}, ``Multi-modal data acquisition platform for behavioral evaluation,'' Harvard Dataverse, 2024. DOI: 10.7910/DVN/HMZ5RG.

\bibitem{akiba2019}
T.~Akiba, S.~Sano, T.~Yanase, T.~Ohta, and M.~Koyama, ``Optuna: A next-generation hyperparameter optimization framework,'' in \emph{Proc. ACM SIGKDD Int. Conf. Knowl. Discov. Data Mining}, 2019, pp.~2623--2631.

\bibitem{breiman2001}
L.~Breiman, ``Random forests,'' \emph{Mach. Learn.}, vol.~45, no.~1, pp.~5--32, 2001.

\bibitem{akerstedt1990}
T.~\r{A}kerstedt and M.~Gillberg, ``Subjective and objective sleepiness in the active individual,'' \emph{Int. J. Neurosci.}, vol.~52, no.~1--2, pp.~29--37, 1990.

\bibitem{platt1999}
J.~C. Platt, ``Probabilistic outputs for support vector machines and comparisons to regularized likelihood methods,'' in \emph{Advances in Large Margin Classifiers}, MIT Press, 1999, pp.~61--74.

\bibitem{han2005}
H.~Han, W.-Y. Wang, and B.-H. Mao, ``Borderline-SMOTE: A new over-sampling method in imbalanced data sets learning,'' in \emph{Proc. Int. Conf. Intell. Comput.}, Springer, 2005, pp.~878--887.

\bibitem{he2008adasyn}
H.~He, Y.~Bai, E.~A. Garcia, and S.~Li, ``ADASYN: Adaptive synthetic sampling approach for imbalanced learning,'' in \emph{Proc. IEEE Int. Joint Conf. Neural Netw.}, 2008, pp.~1322--1328.

\bibitem{chen2016}
T.~Chen and C.~Guestrin, ``XGBoost: A scalable tree boosting system,'' in \emph{Proc. ACM SIGKDD Int. Conf. Knowl. Discov. Data Mining}, 2016, pp.~785--794.

\bibitem{bandt2002}
C.~Bandt and B.~Pompe, ``Permutation entropy: A natural complexity measure for time series,'' \emph{Phys. Rev. Lett.}, vol.~88, no.~17, p.~174102, 2002.

\bibitem{shrout1979}
P.~E. Shrout and J.~L. Fleiss, ``Intraclass correlations: Uses in assessing rater reliability,'' \emph{Psychol. Bull.}, vol.~86, no.~2, pp.~420--428, 1979.

\end{thebibliography}

\end{document}
